%------------------------------------------------------------------------------%
\section{Transformação de Funções}

Nessa seção usaremos translações, expansões e reflexões para fazer o esboço de
gráfico de algumas funções a partir de funções já conhecidas (seção anterior).
Abaixo, listamos algumas transformações nas variáveis $x$ e $y$ e como elas
afetam o gráfico de um função $y=f(x)$.

%------------------------------------------------------------------------------%
\begin{center}
\begin{tabular}{l<{\quad}>{\quad}l}
\hline
  $f(x) \mapsto  f(-x)$  & Reflexão em torno do eixo $y$                 \\
  $f(x) \mapsto -f(x)$   & Reflexão em torno do eixo $x$                 \\
  $f(x) \mapsto  f(x+a)$ & Translação horizontal                         \\
                         & direita se $a < 0$, esquerda se $a > 0$       \\
  $f(x) \mapsto f(x)+a$  & Translação vertical,                          \\
                         & abaixo se $a < 0$, acima se $a > 0$           \\
  $f(x) \mapsto f(ax)$   & Mudança de escala horizontal                  \\
                         & contração se $a > 1$, expansão se $0 < a < 1$ \\
  $f(x) \mapsto af(x)$   & Mudança de escala vertical, $a>0$             \\
                         & expansão se $a > 1$, contração se $0 < a < 1$ \\
\hline
\end{tabular}
\end{center}
%------------------------------------------------------------------------------%

%------------------------------------------------------------------------------%
\begin{center}
\begin{tabular}{ccc}
Original & Reflexão em $y$ & Reflexão em $x$  \\
\includegraphics[scale=0.25]{geral} &
\includegraphics[scale=0.25]{reflecty} &
\includegraphics[scale=0.25]{reflectx} \\
$y=f(x)$ & $y=f(-x)$ & $y=-f(x)$ \\
\end{tabular}
\end{center}
%------------------------------------------------------------------------------%

%------------------------------------------------------------------------------%
\begin{center}
\begin{tabular}{cccc}
Translação direita &
Translação esquerda &
Translação acima &
Translação abaixo \\
\includegraphics[scale=0.25]{transh} &
\includegraphics[scale=0.25]{transh2} &
\includegraphics[scale=0.25]{transv} &
\includegraphics[scale=0.25]{transv2}\\
$y=f(x+a), a< 0$ & $y=f(x+a), a> 0$ & $y=f(x)+a, a>0$ & $y=f(x)+a, a<0$ \\
\end{tabular}
\end{center}
%------------------------------------------------------------------------------%

%------------------------------------------------------------------------------%
\begin{center}
\begin{tabular}{cccc}
Expansão horizontal &
Contração horizontal &
Expansão vertical &
Contração vertical \\
\includegraphics[scale=0.25]{scaleh} &
\includegraphics[scale=0.25]{scaleh2} &
\includegraphics[scale=0.25]{scalev} &
\includegraphics[scale=0.25]{scalev2} \\
$y=f(ax), 0<a< 1$ & $y=f(ax), a> 1$ & $y=af(x), a>1$ & $y=af(x), 0<a<1$  \\
\end{tabular}
\end{center}
%------------------------------------------------------------------------------%

%------------------------------------------------------------------------------%
\begin{example}
  A partir do gráfico de $f(x)= \sqrt{x}$,
  use transformações para obter os gráficos de
  \begin{mathlist}
    \item y =  \sqrt{x}-1
    \item y =  \sqrt{x-2}
    \item y = -\sqrt{x}
    \item y = 2\sqrt{x}
    \item y =  \sqrt{-x}
  \end{mathlist}
\end{example}
%------------------------------------------------------------------------------%
